% GENERUOJAMA is rezultatai/darbiniai/sprendimu_matrica.csv
% Ranka NELIESTI - paleisti: python -m src.eksperimentai.matrica
\begingroup
\footnotesize
\setlength{\tabcolsep}{4pt}
\begin{tabularx}{\textwidth}{@{}>{\raggedright\arraybackslash}X>{\centering\arraybackslash}p{1.15cm}>{\centering\arraybackslash}p{1.35cm}>{\centering\arraybackslash}p{1.15cm}>{\centering\arraybackslash}p{1.9cm}>{\centering\arraybackslash}p{1.4cm}>{\raggedright\arraybackslash}p{1.4cm}@{}}
\toprule
\textbf{Metodas} & \textbf{Kokybė} & \textbf{Disba\-lansas} & \textbf{Resur\-sai} & \textbf{Interpre\-tuo\-jamumas} & \textbf{Svertinė suma} & \textbf{Šalt.} \\
\multicolumn{1}{@{}l}{\itshape svoris} & 0,30 & 0,30 & 0,25 & 0,15 & --- & \\
\midrule
\multicolumn{7}{@{}l}{\textit{Prižiūrimi (8 kategorijų klasifikavimas)}} \\
\addlinespace[2pt]
XGBoost & 5 & 5 & 5 & 3 & \textbf{4,70} & \cite{almahaqeri2026gradient} \\
LightGBM & 5 & 5 & 5 & 3 & \textbf{4,70} & \cite{almahaqeri2026gradient} \\
Sprendimų medis (CART) & 3 & 3 & 5 & 5 & \textbf{3,80} & (aut.) \\
Random Forest & 4 & 3 & 4 & 3 & \textbf{3,55} & \cite{imani2025imbalance} \\
Daugiasluoksnis perceptronas & 4 & 3 & 4 & 1 & \textbf{3,25} & \cite{mazinani2026constrained} \\
1D-CNN & 3 & 3 & 3 & 1 & \textbf{2,70} & (aut.) \\
\addlinespace
\multicolumn{7}{@{}l}{\textit{Neprižiūrimi (dvejetainė anomalijų formuluotė)}} \\
\addlinespace[2pt]
Isolation Forest & 2 & 4 & 5 & 3 & \textbf{3,50} & (aut.) \\
Autokoderis & 2 & 4 & 4 & 1 & \textbf{2,95} & \cite{meidan2018nbaiot} \\
\bottomrule
\end{tabularx}
\endgroup
